\cotefile{109}
\scanpage{13}
\begin{center}\uline{COMMENTS ON ``PURSUING STACKS''}\end{center}

\diagram{\begin{tabular}{@{}l@{\quad}l@{}}
p2 $\ell$8 & inadequate\\
p2-$\ell$-15 & incomplete\\
p3 $\ell$-8 & pertinent\\
p4 $\ell$6 & indiscriminate\\
\hspace*{1.45em}$\ell$-1 & misled\\
p4-$\ell$-16 & arbitrary\\
\hspace*{1.45em}$\ell$-8 & the same\\
\multicolumn{2}{@{}l@{}}{Section 14 etc}\\
p2 $\ell$18 & impediment\\
p3 $\ell$4 & CW-spaces\\
\hspace*{1.45em}$\ell$12 & have\\
\hspace*{1.45em}$\ell$15 & adequate
\end{tabular}}

p3\qquad The question of cubical versus simplicial sets is an interesting one. The category of simplicial sets has a number of good properties, which have made it very popular.

1.\quad It is cartesian closed.

2.\quad The realisation functor gives an equivalence between the homotopy category of Kan simplicial sets and the homotopy category of CW-spaces and the proof of this has been published.

3.\quad The realisation functor satisfies $\lvert X\times Y\rvert$ is homeomorphic to $\lvert X\rvert\times\lvert Y\rvert$ (at least if one takes the fine topology with respect to compact subsets on the latter product, as is the correct thing to do for other purposes, such as obtaining the right cartesian closed category for topological spaces). Also simplicial sets are firmly embedded in the literature. Again, there is a Dold-Kan theorem, giving an equivalence between chain complexes and simplicial abelian groups. A neat proof of the Serre spectral sequence using simplicial methods has been published by A. Dress (see the reference in my paper with Heath ``Cogluing homotopy equivalences'').

By contrast, cubical sets are not greatly used in the literature, although a number of writers have found them convenient. A recent exposition by Massey of singular homology theory uses the cubical method throughout, and finds them very convenient. Again, Brown-Higgins have found them essential for their methods. The difficulties about using them are that they are not firmly embedded in the
% archive sheet number: 12
